% ==============================================================================
% PhD Curriculum Vitae — High-Density Template
% ==============================================================================
\documentclass[10pt, letterpaper]{article}

%------------------------------------------------------------------------------
% Packages
%------------------------------------------------------------------------------
\usepackage[margin=0.75in, top=0.65in, bottom=0.65in]{geometry}
\usepackage{ctex} % 中文支持（关键）
\usepackage{xeCJK}
\setCJKmainfont{Noto Serif CJK SC}
\usepackage{hyperref}
\usepackage{fontawesome5}
\usepackage{titlesec}
\usepackage{enumitem}
\usepackage{tabularx}
\usepackage{array}
\usepackage{xcolor}
\usepackage{microtype}
\usepackage[T1]{fontenc}
\usepackage{lmodern}
\usepackage{ifthen}
\definecolor{RULECOLOR}{RGB}{180,180,180}

%------------------------------------------------------------------------------
% Colors
%------------------------------------------------------------------------------
\definecolor{accentblue}{RGB}{20, 75, 140}
\definecolor{darkgray}{RGB}{40, 40, 40}
\definecolor{lightgray}{RGB}{120, 120, 120}
\definecolor{rulecolor}{RGB}{180, 200, 225}

%------------------------------------------------------------------------------
% Hyperref
%------------------------------------------------------------------------------
\hypersetup{colorlinks=true, urlcolor=accentblue, linkcolor=accentblue}

%------------------------------------------------------------------------------
% Spacing
%------------------------------------------------------------------------------
\setlength{\parskip}{0pt}
\setlength{\parindent}{0pt}

%------------------------------------------------------------------------------
% Section formatting
%------------------------------------------------------------------------------
\titleformat{\section}
  {\normalsize\bfseries\color{accentblue}\uppercase}
  {}{0em}{}
  [\vspace{-5pt}\textcolor{rulecolor}{\rule{\linewidth}{0.5pt}}\vspace{1pt}]
\titlespacing{\section}{0pt}{8pt}{3pt}

\titleformat{\subsection}
  {\small\bfseries\color{darkgray}}
  {}{0em}{}
\titlespacing{\subsection}{0pt}{4pt}{1pt}

%------------------------------------------------------------------------------
% List settings
%------------------------------------------------------------------------------
\setlist{noitemsep, topsep=1pt, parsep=0pt}
\setlist[itemize]{leftmargin=1.2em, label=\textcolor{accentblue}{\textbullet}}

%------------------------------------------------------------------------------
% Custom commands
%------------------------------------------------------------------------------
\newcommand{\entry}[3]{%
  \begin{tabularx}{\linewidth}{@{}X r@{}}
    \textbf{#1} \textcolor{lightgray}{$\cdot$} \textit{#2} &
    {\small\textcolor{lightgray}{#3}}
  \end{tabularx}\vspace{1pt}%
}

\newcommand{\entrysub}[4]{%
  \begin{tabularx}{\linewidth}{@{}X r@{}}
    \textbf{#1} \textcolor{lightgray}{$\cdot$} \textit{#2} &
    {\small\textcolor{lightgray}{#3}}
  \end{tabularx}%
  \ifthenelse{\equal{#4}{}}{}{\\\hspace{1.2em}{\small\textcolor{darkgray}{#4}}}%
  \vspace{1pt}%
}

\newcommand{\honorrow}[3]{%
  \begin{tabularx}{\linewidth}{@{}X r@{}}
    #1 \textcolor{lightgray}{$\cdot$} \textit{\small #2} & {\small\textcolor{lightgray}{#3}}
  \end{tabularx}\vspace{1pt}%
}

%==============================================================================
\begin{document}
\pagestyle{empty}

%==============================================================================
% HEADER
%==============================================================================
\begin{center}
  {\Large\bfseries\color{darkgray} 文曲星}\\[3pt]
  {\small\color{lightgray}
    博士生 $\cdot$ 计算机科学与技术学院 $\cdot$ 清华大学\\[2pt]
    \faEnvelope[regular]\;\href{mailto:name@qinghua.edu.cn}{mail@tsinghua.edu.cn}\quad
    \faGraduationCap\;\href{https://scholar.google.com}{Google Scholar}\quad
    \faMobile*\;+86 123-4567-8910\quad
    \faMapMarker*\;中国北京
  }
\end{center}

%==============================================================================
\section{研究兴趣}
机器人感知与自主导航方向，融合深度学习与多传感器融合技术，围绕视觉SLAM、三维场景理解与语义地图构建开展研究，重点探索视觉--惯性里程计、动态环境下的鲁棒定位以及多机器人协同建图框架，推动高效、可泛化的智能移动机器人系统发展。

%==============================================================================
\section{教育经历}

\entrysub{清华大学（Tsinghua University）}{博士在读-计算机科学与技术专业}{2023年9月--2027年6月（预计）}%
{导师：博士导师\quad 中国北京}

\entrysub{哈尔滨工业大学（Harbin Institute of Technology）}{工学硕士-控制科学与工程专业}{2020年9月--2023年6月}%
{导师：硕士导师\quad 中国哈尔滨}

\entrysub{苏黎世联邦理工学院（ETH Zurich）}{访问学者-计算机视觉专业}{2022年3月--2022年12月}%
{导师：Prof.\ Marc Pollefeys\quad 瑞士苏黎世}

\entrysub{武汉大学（Wuhan University）}{工学学士-自动化专业}{2016年9月--2020年6月}%
{中国武汉}

%==============================================================================
\section{发表论文}
{\small\textcolor{lightgray}{*Corresponding author}}
{\faExternalLink*\;\href{https://scholar.google.com}{Full list on Google Scholar.}}
\vspace{2pt}

\begin{enumerate}[label={\small[\arabic*]}, itemsep=3pt, topsep=0pt]

\item \textbf{Wen, Q.}, Chen H, Liu Y, et al.
    ``Tile: Paper Title.''
    \textit{Proc.\ RSS}, 2026.
    [\href{https://arxiv.org/}{arXiv}][\href{https://github.com/}{Code}]

\item \textbf{Wen, Q.}, Chen H, Liu Y, et al.
    ``Tile: Paper Title.''[C]
    \textit{Proceedings of the IEEE/RSJ International Conference on Intelligent Robots and Systems (IROS). 2099: 2210-2218.}
    [\href{https://github.com/}{Code}]
    [\href{https://}{Demo}]

\item Zhang W, \textbf{Wen, Q.}, Liu H, et al.
    ``Tile: Paper Title.''[J]
    \textit{IEEE Transactions on Intelligent Transportation Systems, 2099, 26(3): 1520-1538.}

\item Chen, H, \textbf{Wen, Q.}, Wu Y, et al.
    ``Tile: Paper Title.''[C]
   \textit{Proc.\ IEEE ICRA 2025}.

\end{enumerate}
%==============================================================================
\section{研究经历}


\entry{Title: 论文标题}{清华大学智能机器人实验室}{2024年9月 -- 2026年3月}
\begin{itemize}
\item 设计基于神经隐式表征的多机器人协同稠密建图框架
\item 构建分布式地图融合架构，支持5台机器人实时协同
\item 在3个大规模室内外数据集上建图精度超越SOTA方法12.5\%
\end{itemize}

\entry{Title: 论文标题}{清华大学智能机器人实验室}{2023年9月 -- 2024年8月}
\begin{itemize}
\item 构建动态物体感知的视觉SLAM系统
\item 融合语义分割与光流估计进行动态区域剔除
\item 在TUM RGB-D动态序列上定位精度提升38\%
\end{itemize}

\entry{Title: 论文标题}{ETH Zurich}{2022年3月 -- 2022年10月}
\begin{itemize}
\item 开发基于深度学习的单目视觉惯性里程计系统
\item 构建包含50K帧的室内外导航数据集
\end{itemize}

\entry{Title: 论文标题}{哈尔滨工业大学}{2021年9月 -- 2021年12月}
\begin{itemize}
\item 华为杯第18届全国研究生数学建模竞赛
\item 基于ICP与NDT算法的点云配准优化
\item 实现0.12m平均配准误差
\end{itemize}
%==============================================================================
\section{业界经历}

\entry{算法实习生（Research Intern）}{大疆创新（DJI）感知算法部，中国深圳}{2099年3月--2099年8月}

\begin{itemize}
\item 参与无人机视觉SLAM模块优化，定位精度提升15\%
\item 开发基于点云的障碍物检测算法，部署于嵌入式平台
\end{itemize}

\entry{研究助理（Research Assistant）}{苏黎世联邦理工学院（ETH Zurich）计算机视觉实验室，瑞士}{2099年8月--2099年12月}

\entry{机器人开发实习生（Robotics Intern）}{海康威视机器人事业部，中国杭州}{2099年6月--2099年9月}

\entry{软件开发实习生（Software Engineering Intern）}{华为技术有限公司，中国深圳}{2099年7月--2099年10月}
%==============================================================================
\section{荣誉奖项}
\honorrow{电气与电子工程师协会会员（IEEE Membership）}{IEEE}{2024--2027}
\honorrow{中国自动化学会会员（CAA Membership）}{CAA}{2021--2025}
\honorrow{省级一等奖}{湖北省数学建模竞赛}{2018}
\honorrow{优秀一等奖学金}{武汉大学}{2017--2020}

%==============================================================================
\section{其他}

\begin{tabularx}{\linewidth}{@{}>{\bfseries\small}p{2.5cm}X@{}}
编程语言 & Python, C++, CUDA, Bash, SQL, ROS \\

技术框架 & PyTorch, TensorRT, Open3D, OpenCV, PCL, ORB-SLAM3, GTSAM \\[1.5pt]

基础设施 & Git, Docker/Singularity, Slurm, AWS（EC2/S3）, Weights \& Biases \\[1.5pt]

语言能力 & 中文（母语）, 英语（流利，IELTS 7.5）, 德语（初级，A2） \\[1.5pt]

兴趣爱好 & 无人机航拍, 骑行, 乒乓球, 科技博客写作 \\

推荐人 & 院士教授（导师, yuanshi@tsinghua.edu.cn） \\
\end{tabularx}

\vspace{8pt}
\noindent\textcolor{lightgray}{\footnotesize\textit{Last updated: \today}}

\end{document}